\Needspace{4\baselineskip}
\begin{Example}
\label{ex:differentiation-operator-closure-sobolev}
现在回到 Example~\ref{ex:differentiation-operator-domain-and-unboundedness} 中定义的微分算子 $(\mathcal{D}(D),D)$. 在上述评论中, 提到该算子的定义域并不 ``最优'', 因为它未包括分片可微函数这样显然应考虑的函数. 现在可以看出, 问题在于该算子不是闭的. 按照上面 Lemma~\ref{lem:closable-operator-zero-sequence-criterion} 的证明中的方法构造闭包. 即将闭包的定义域 $\mathcal{D}(\overline D)$ 定义为满足以下条件的函数 $u\in L^2(0,L)$: 存在序列 $u_n\in\mathcal{D}(D)$ 及元素 $v\in L^2(0,L)$, 使得\begin{enumerate}\item $u_n\to u$ 于 $L^2(0,L)$ 中, 且\item $\dfrac{\mathrm{d}u_n}{\mathrm{d}x}\to v$ 于 $L^2(0,L)$ 中.\end{enumerate}然而, 这正是 Sobolev 空间的定义:\[\mathcal{D}(\overline D):=\{u\in H^1(0,L)\mid u(0)=0\},\]其中的边界条件在迹的意义下理解.
\end{Example}

\Needspace{4\baselineskip}
\begin{Example}
\label{ex:laplacian-dirichlet-closure-partial-proof}
以大致相同的方法可使 Laplacian 算子闭合. 考察 Dirichlet 条件的问题. 假设存在 $\widetilde{\mathcal{D}}_D(\Delta)$ 中的序列 $u_n$ 及 $L^2(\Omega)$ 中的函数 $u$ 和 $v$, 使得
\begin{equation}\label{eq:laplacian-closure-function-convergence}
u_n\to u\text{ in }L^2(\Omega),
\end{equation}
以及
\begin{equation}\label{eq:laplacian-closure-image-convergence}
\Delta u_n\to v\text{ in }L^2(\Omega).
\end{equation}
将在 Chapter~\ref{chap:linear-elliptic-equations} 中证明, 扩张后的定义域为 $\mathcal{D}_D(\Delta)=H^2(\Omega)\cap H_0^1(\Omega)$. 虽然目前还无法完全论证这一断言, 但可以给出部分证明. 首先注意, 若 $\Omega$ 为 $C^1$ 类, 则有 $\mathcal{D}_D(\Delta)\supset H^2(\Omega)\cap H_0^1(\Omega)$. 为看出这一点, 注意 $\mathcal{D}_D(\Delta)$ 是 $\widetilde{\mathcal{D}}_D(\Delta)$ 关于 $H^2(\Omega)$ 范数的完备化的子集. 还要注意, 根据 Theorem~\ref{thm:zero-boundary-sobolev-characterization}, $H^2(\Omega)\cap H_0^1(\Omega)$ 是在迹的意义下满足边界 $\partial\Omega$ 上 $u=0$ 的所有 $u\in H^2(\Omega)$ 所组成的集合. 然而, 再次利用迹定理可知, 这个集合也是 $\widetilde{\mathcal{D}}_D(\Delta)$ 关于 $H^2(\Omega)$ 范数的完备化. 其次注意, $\mathcal{D}_D(\Delta)\subset H^1(\Omega)$. 为看出这一点, 利用 Green 恒等式得到
\begin{equation}\label{eq:laplacian-closure-green-energy-identity}
\int_\Omega|\nabla u_n|^2=-\int_\Omega u_n\Delta u_n.
\end{equation}
于是应用 Hölder 不等式可知, 满足 1 和 2 的序列 $u_n\in\mathcal{D}_D(\Delta)$ 在 $H^1(\Omega)$ 中有界. 因而根据 Theorem~\ref{thm:alaoglu-weak-star-subsequence} (以及 $H^1(\Omega)$ 的自反性), $u_n$ 有在 $H^1(\Omega)$ 中弱收敛的子序列 $u_n\rightharpoonup\overline u$. 然而, 根据弱极限的唯一性, 必有 $u=\overline u$.
\end{Example}

\Needspace{4\baselineskip}
\begin{Example}
\label{ex:differential-operators-closable-wave-domain}
注意, 所有具有光滑系数的微分算子都是可闭的. 为看出这一点, 假设 $(\mathcal{D}(L),L)$ 是一个 $m$ 阶微分算子, 且 $\mathcal{D}(L)\subset L^2(\Omega)$. 那么若 $x_n\to0$ 于 $L^2(\Omega)$ 中, 就有 $L(x_n)\to0$ 于 $H^{-m}(\Omega)$ 中. 由于 $L^2(\Omega)$ 连续嵌入 $H^{-m}(\Omega)$, 所以 $\lim L(x_n)$ (看作 $L^2(\Omega)$ 中的极限) 要么等于 $0$, 要么不存在. 因而根据 Lemma~\ref{lem:closable-operator-zero-sequence-criterion}, 该算子可闭.

虽然所有微分算子都是可闭的, 但闭包的定义域未必如人们所预期的那样是 Sobolev 空间. 例如, 考察波动算子
\begin{equation}\label{eq:wave-operator-closable-nonsobolev-example}
L(u):=u_{tt}-u_{xx},
\end{equation}
其定义域为 $H^2(\mathbb{R}\times(0,T))$, 其中 $T<\infty$. 请读者验证, 形如 $u(x,t)=f(x-t)$ 的显然的波形解属于该算子的闭包, 其中 $f\in L^2(\mathbb{R})$ 不连续. 当然, 这样的函数既不属于 $H^2(\mathbb{R}\times(0,T))$, 也不属于 $H^1(\mathbb{R}\times(0,T))$. (它们沿不连续线没有良好定义的迹.)
\end{Example}

\Needspace{4\baselineskip}
\begin{Example}
\label{ex:nonclosable-integral-operator}
虽然实际遇到的大多数算子都是可闭的, 但也存在不可闭的算子. 设 $X=L^2(\mathbb{R})$, $Y=\mathbb{R}$, 令 $\mathcal{D}(A)$ 为具有紧支集的有界函数的集合, 并令 $Au:=\int_{-\infty}^\infty u(x)\,\mathrm{d}x$, 以此定义 $(\mathcal{D}(A),A)$. 那么若取
\begin{equation}\label{eq:nonclosable-integral-operator-sequence}
u_n(x)=\begin{cases}\dfrac{1}{n},&|x|\leq n,\\0,&|x|>n,\end{cases}
\end{equation}
就可看出, $u_n\to0$ 于 $L^2(\mathbb{R})$ 中, 但 $Au_n=2$.
\end{Example}
